\PassOptionsToPackage{unicode}{hyperref}
\documentclass[11pt,a4paper]{article}
\usepackage[T1]{fontenc}
\usepackage[utf8]{inputenc}
\usepackage{lmodern,microtype,xcolor}
\usepackage{tcolorbox}
\tcbuselibrary{skins,breakable,listings}
\usepackage{titlesec,fancyhdr,enumitem,booktabs,longtable,array,colortbl}
\usepackage[margin=2.4cm]{geometry}
\usepackage{listings}
\usepackage{hyperref}

\definecolor{javaOrange}{HTML}{E76F00}
\definecolor{javaDeep}{HTML}{BF5700}
\definecolor{javaBlue}{HTML}{1565C0}
\definecolor{javaTeal}{HTML}{00695C}
\definecolor{javaAmber}{HTML}{B45309}
\definecolor{javaInk}{HTML}{2B2140}
\definecolor{javaCodeBg}{HTML}{FFF8F0}
\definecolor{javaCodeInline}{HTML}{BF5700}
\definecolor{javaRule}{HTML}{FFD7B5}
\definecolor{javaQuoteBg}{HTML}{FFFAF5}
\color{javaInk}

\lstdefinestyle{java}{
  language=Java,
  basicstyle=\small\ttfamily\color{javaInk},
  keywordstyle=\color[HTML]{D73A49}\bfseries,
  stringstyle=\color[HTML]{032F62},
  commentstyle=\color[HTML]{6A737D}\itshape,
  breaklines=true,breakatwhitespace=true,
  tabsize=4,showstringspaces=false,
  morekeywords={var,record,String,Scanner,ArrayList,HashMap,System,Override,
    public,private,protected,class,interface,extends,implements,new,this,super,
    return,void,static,final,abstract,null,true,false,throws,throw,try,catch,finally},
}
\newtcblisting{javacode}{
  listing only,listing style=java,
  breakable,enhanced jigsaw,arc=4pt,boxrule=0pt,frame hidden,
  colback=javaCodeBg,borderline west={3pt}{0pt}{javaDeep},
  left=10pt,right=8pt,top=6pt,bottom=6pt,
  before skip=8pt,after skip=8pt,
}
\renewcommand{\texttt}[1]{{\color{javaCodeInline}\small\ttfamily #1}}
\renewenvironment{quote}
  {\begin{tcolorbox}[breakable,enhanced jigsaw,arc=4pt,boxrule=0pt,frame hidden,
    colback=javaQuoteBg,borderline west={3pt}{0pt}{javaOrange},
    left=12pt,right=10pt,top=8pt,bottom=8pt,before skip=8pt,after skip=8pt]\itshape}
  {\end{tcolorbox}}
\titleformat{\section}{\sffamily\Large\bfseries\color{javaOrange}}{\thesection}{0.7em}{}
  [{\vspace{2pt}\color{javaRule}\titlerule[1.2pt]}]
\titleformat{\subsection}{\sffamily\large\bfseries\color{javaDeep}}{\thesubsection}{0.6em}{}
\titleformat{\subsubsection}{\sffamily\normalsize\bfseries\color{javaTeal}}{\thesubsubsection}{0.5em}{}
\titleformat{\paragraph}[runin]{\sffamily\bfseries\color{javaAmber}}{}{0em}{}
\titlespacing*{\section}{0pt}{2.2ex plus 1ex minus .2ex}{1.2ex plus .2ex}
\setlength{\headheight}{24pt}
\pagestyle{fancy}\fancyhf{}
\renewcommand{\headrulewidth}{0.6pt}
\fancyhead[L]{\sffamily\small\color{javaDeep}Java Programming MOOC}
\fancyhead[R]{\sffamily\small\color{javaDeep}\leftmark}
\fancyfoot[C]{\sffamily\small\color{javaOrange}\thepage}
\renewcommand{\headrule}{\hbox to\headwidth{\color{javaRule}\leaders\hrule height \headrulewidth\hfill}}
\setlist[itemize]{leftmargin=1.4em,itemsep=2pt,topsep=3pt}
\setlist[enumerate]{leftmargin=1.6em,itemsep=2pt,topsep=3pt}
\renewcommand{\labelitemi}{\textcolor{javaOrange}{\textbullet}}
\renewcommand{\arraystretch}{1.25}
\arrayrulecolor{javaRule}
\hypersetup{colorlinks=true,linkcolor=javaDeep,urlcolor=javaBlue,citecolor=javaTeal}
\makeatletter
\newcommand{\subtitle}[1]{\gdef\@subtitle{#1}}
\providecommand{\@subtitle}{}
\def\@maketitle{%
  \newpage\null\vskip 2em%
  \begin{center}%
    {\sffamily\bfseries\Huge\color{javaOrange}\@title\par}\vskip 1em%
    {\sffamily\Large\color{javaDeep}\@subtitle\par}\vskip 1.2em%
    {\color{javaRule}\rule{0.6\linewidth}{1.4pt}\par}\vskip 1em%
    {\sffamily\large\color{javaInk}\@author\par}\vskip 0.4em%
    {\sffamily\small\color{javaAmber}\@date\par}%
  \end{center}\par\vskip 2em}
\makeatother

\title{Programação em Java}
\subtitle{Lição 6: Listas de Objetos, Testes e Programas Complexos}
\author{Java Programming MOOC · Universidade de Helsinque --- tradução para o português}
\date{2026}

\begin{document}
\maketitle
{\hypersetup{linkcolor=javaDeep}\setcounter{tocdepth}{3}\tableofcontents}
\newpage

\section{Lição 6: Listas de Objetos, Testes e Programas Complexos}

\subsection{Objetivos de aprendizagem}

Ao final desta lição, você será capaz de:

\begin{itemize}
\item Usar listas de objetos e objetos que contêm listas
\item Separar a interface do usuário da lógica do programa
\item Escrever e executar testes unitários com JUnit
\item Organizar programas maiores em múltiplas classes
\item Entender o conceito de desenvolvimento orientado a testes (TDD)
\end{itemize}

\subsection{Objetos em listas e listas dentro de objetos}

\subsubsection{Lista de objetos}

Já vimos que podemos armazenar objetos em listas. Vejamos um exemplo mais completo:

\begin{javacode}
ArrayList<Livro> biblioteca = new ArrayList<>();
biblioteca.add(new Livro("Dom Casmurro", "Machado de Assis", 1899));
biblioteca.add(new Livro("O Cortico", "Aluisio Azevedo", 1890));
biblioteca.add(new Livro("Vidas Secas", "Graciliano Ramos", 1938));
\end{javacode}

Buscando um livro por título:

\begin{javacode}
public static Livro buscarPorTitulo(ArrayList<Livro> lista, String titulo) {
    for (Livro livro : lista) {
        if (livro.getTitulo().equalsIgnoreCase(titulo)) {
            return livro;
        }
    }
    return null;
}
\end{javacode}

Filtrando por ano:

\begin{javacode}
public static ArrayList<Livro> livrosAntesDe(ArrayList<Livro> lista, int ano) {
    ArrayList<Livro> resultado = new ArrayList<>();
    for (Livro livro : lista) {
        if (livro.getAno() < ano) {
            resultado.add(livro);
        }
    }
    return resultado;
}
\end{javacode}

\subsubsection{Objeto com lista interna}

Uma classe pode ter um \texttt{ArrayList} como variável de instância:

\begin{javacode}
public class Turma {
    private String nome;
    private ArrayList<Aluno> alunos;

    public Turma(String nome) {
        this.nome = nome;
        this.alunos = new ArrayList<>();
    }

    public void adicionarAluno(Aluno aluno) {
        this.alunos.add(aluno);
    }

    public int quantidadeDeAlunos() {
        return this.alunos.size();
    }

    public double mediaGeral() {
        if (this.alunos.isEmpty()) return 0.0;

        double somaMedias = 0;
        for (Aluno aluno : this.alunos) {
            somaMedias += aluno.calcularMedia();
        }
        return somaMedias / this.alunos.size();
    }

    public void listarAprovados(double notaMinima) {
        System.out.println("Aprovados em " + this.nome + ":");
        for (Aluno aluno : this.alunos) {
            if (aluno.calcularMedia() >= notaMinima) {
                System.out.println("  " + aluno);
            }
        }
    }

    @Override
    public String toString() {
        return this.nome + " (" + this.alunos.size() + " alunos)";
    }
}
\end{javacode}

\subsection{Separando a interface do usuário da lógica do programa}

Um princípio fundamental de bom design é \textbf{separar} o que o programa faz (lógica) de como ele interage com o usuário (interface).

\subsubsection{O problema de misturar tudo}

\begin{javacode}
// Ruim: logica e UI juntas
public static void main(String[] args) {
    Scanner scanner = new Scanner(System.in);
    ArrayList<String> itens = new ArrayList<>();

    System.out.print("Quantos itens? ");
    int n = Integer.valueOf(scanner.nextLine());

    for (int i = 0; i < n; i++) {
        System.out.print("Item " + (i+1) + ": ");
        itens.add(scanner.nextLine());
    }
    // logica misturada com UI...
}
\end{javacode}

\subsubsection{A solução: classes separadas}

\textbf{Classe de lógica:}

\begin{javacode}
public class ListaDeCompras {
    private ArrayList<String> itens;

    public ListaDeCompras() {
        this.itens = new ArrayList<>();
    }

    public void adicionar(String item) {
        this.itens.add(item);
    }

    public boolean remover(String item) {
        return this.itens.remove(item);
    }

    public boolean contem(String item) {
        return this.itens.contains(item);
    }

    public int tamanho() {
        return this.itens.size();
    }

    public ArrayList<String> getItens() {
        return new ArrayList<>(this.itens); // retorna copia
    }

    @Override
    public String toString() {
        return this.itens.toString();
    }
}
\end{javacode}

\textbf{Classe de interface:}

\begin{javacode}
public class InterfaceListaDeCompras {
    private Scanner scanner;
    private ListaDeCompras lista;

    public InterfaceListaDeCompras(Scanner scanner) {
        this.scanner = scanner;
        this.lista = new ListaDeCompras();
    }

    public void executar() {
        while (true) {
            System.out.println("\n1. Adicionar  2. Remover  3. Listar  4. Sair");
            System.out.print("Escolha: ");
            String opcao = scanner.nextLine();

            if (opcao.equals("1")) {
                adicionarItem();
            } else if (opcao.equals("2")) {
                removerItem();
            } else if (opcao.equals("3")) {
                listarItens();
            } else if (opcao.equals("4")) {
                System.out.println("Ate logo!");
                break;
            } else {
                System.out.println("Opcao invalida.");
            }
        }
    }

    private void adicionarItem() {
        System.out.print("Item: ");
        String item = scanner.nextLine();
        this.lista.adicionar(item);
        System.out.println("'" + item + "' adicionado.");
    }

    private void removerItem() {
        System.out.print("Item a remover: ");
        String item = scanner.nextLine();
        if (this.lista.remover(item)) {
            System.out.println("Removido.");
        } else {
            System.out.println("Item nao encontrado.");
        }
    }

    private void listarItens() {
        System.out.println("Sua lista (" + this.lista.tamanho() + " itens):");
        for (String item : this.lista.getItens()) {
            System.out.println("  - " + item);
        }
    }
}
\end{javacode}

\textbf{\texttt{main} simples:}

\begin{javacode}
public static void main(String[] args) {
    Scanner scanner = new Scanner(System.in);
    InterfaceListaDeCompras ui = new InterfaceListaDeCompras(scanner);
    ui.executar();
}
\end{javacode}

\subsection{Introdução a testes unitários}

\textbf{Testes unitários} verificam automaticamente se um método ou classe se comporta corretamente. Em Java, usamos a biblioteca \textbf{JUnit}.

\subsubsection{Por que testar?}

\begin{itemize}
\item Detecta bugs antes de chegar ao usuário
\item Dá confiança para refatorar o código
\item Serve como documentação do comportamento esperado
\end{itemize}

\subsubsection{Estrutura de um teste JUnit}

\begin{javacode}
import org.junit.Test;
import static org.junit.Assert.*;

public class CalculadoraTest {

    @Test
    public void testSoma() {
        Calculadora calc = new Calculadora();
        assertEquals(5, calc.soma(2, 3));
    }

    @Test
    public void testSubtracao() {
        Calculadora calc = new Calculadora();
        assertEquals(1, calc.subtracao(3, 2));
    }

    @Test
    public void testDivisaoPorZero() {
        Calculadora calc = new Calculadora();
        assertEquals(0, calc.divisao(10, 0));
    }
}
\end{javacode}

Classe sendo testada:

\begin{javacode}
public class Calculadora {
    private int valor;

    public Calculadora() {
        this.valor = 0;
    }

    public int soma(int a, int b) {
        return a + b;
    }

    public int subtracao(int a, int b) {
        return a - b;
    }

    public int divisao(int a, int b) {
        if (b == 0) return 0;
        return a / b;
    }
}
\end{javacode}

\subsubsection{Métodos de asserção do JUnit}

\begin{longtable}{ll}
\toprule
\textbf{Método} & \textbf{Verifica} \\
\midrule
\texttt{assertEquals(esperado, real)} & Se dois valores são iguais \\
\texttt{assertNotEquals(naoEsperado, real)} & Se dois valores são diferentes \\
\texttt{assertTrue(condicao)} & Se a condição é verdadeira \\
\texttt{assertFalse(condicao)} & Se a condição é falsa \\
\texttt{assertNull(objeto)} & Se o objeto é \texttt{null} \\
\texttt{assertNotNull(objeto)} & Se o objeto não é \texttt{null} \\
\bottomrule
\end{longtable}

\subsubsection{Desenvolvimento orientado a testes (TDD)}

O \textbf{TDD} (\textit{Test-Driven Development}) inverte o processo tradicional:

\begin{enumerate}
\item \textbf{Escreva o teste} para o comportamento desejado (antes de implementar)
\item \textbf{Execute} --- o teste falha (código ainda não existe)
\item \textbf{Implemente} o mínimo de código para o teste passar
\item \textbf{Execute} --- o teste passa
\item \textbf{Refatore} o código mantendo os testes passando
\end{enumerate}

\begin{javacode}
// Passo 1: escreva o teste primeiro
@Test
public void testListaVaziaRetornaZeroItens() {
    ListaDeCompras lista = new ListaDeCompras();
    assertEquals(0, lista.tamanho());
}

// Passo 1: teste para adicionar item
@Test
public void testAdicionarItemAumentaTamanho() {
    ListaDeCompras lista = new ListaDeCompras();
    lista.adicionar("pao");
    assertEquals(1, lista.tamanho());
}

// Passo 1: teste para remover item
@Test
public void testRemoverItemDiminuiTamanho() {
    ListaDeCompras lista = new ListaDeCompras();
    lista.adicionar("leite");
    lista.remover("leite");
    assertEquals(0, lista.tamanho());
}
\end{javacode}

\subsection{Lendo e interpretando o stack trace}

Quando um erro ocorre, o Java exibe uma trilha de chamadas. Aprenda a ler:

\begin{javacode}
Exception in thread "main" java.lang.NullPointerException
    at Turma.mediaGeral(Turma.java:23)
    at Main.main(Main.java:10)
\end{javacode}

\begin{itemize}
\item \textbf{Tipo do erro}: \texttt{NullPointerException} --- tentativa de usar uma referência \texttt{null}
\item \textbf{Onde}: linha 23 do arquivo \texttt{Turma.java}, dentro do método \texttt{mediaGeral}
\item \textbf{Chamado de}: linha 10 do \texttt{Main.java}
\end{itemize}

Leia de baixo para cima: o \texttt{main} chamou \texttt{mediaGeral}, que falhou na linha 23.

\subsection{Resumo}

\begin{itemize}
\item Listas podem conter objetos; classes podem ter listas como variáveis de instância.
\item \textbf{Separe a lógica da interface do usuário} --- lógica em classes próprias, UI em outra classe.
\item \textbf{Testes unitários} verificam o comportamento de métodos isoladamente.
\item Use \texttt{assertEquals}, \texttt{assertTrue}, \texttt{assertFalse} etc.\ para verificar resultados.
\item \textbf{TDD}: escreva o teste antes do código --- isso guia um design melhor.
\item O \textbf{stack trace} aponta a linha exata onde o erro ocorreu.
\end{itemize}

\subsection{Exercícios}

\paragraph{Exercício 1 --- Biblioteca}

Crie a classe \texttt{Biblioteca} com uma lista de \texttt{Livro}. Implemente métodos: \texttt{adicionarLivro}, \texttt{buscarPorAutor}, \texttt{listarTodosOrdenados} (por título). Crie testes unitários para cada método.

\paragraph{Exercício 2 --- Separação de responsabilidades}

Refatore o seguinte código para separar a lógica da UI:

\begin{javacode}
// Refatore isto:
public static void main(String[] args) {
    Scanner sc = new Scanner(System.in);
    ArrayList<Integer> nums = new ArrayList<>();
    while (true) {
        String l = sc.nextLine();
        if (l.equals("fim")) break;
        nums.add(Integer.valueOf(l));
    }
    int soma = 0;
    for (int n : nums) soma += n;
    System.out.println("Soma: " + soma);
    System.out.println("Media: " + (double) soma / nums.size());
}
\end{javacode}

Crie classes \texttt{ColetorDeNumeros} (lógica) e \texttt{InterfaceColetorDeNumeros} (UI).

\paragraph{Exercício 3 --- TDD: Conta bancária}

Usando TDD, implemente a classe \texttt{ContaBancaria} com os métodos \texttt{depositar}, \texttt{sacar} e \texttt{getSaldo}. Escreva os testes primeiro:
\begin{itemize}
\item \texttt{depositar} aumenta o saldo
\item \texttt{sacar} com saldo suficiente diminui o saldo e retorna \texttt{true}
\item \texttt{sacar} sem saldo suficiente não altera o saldo e retorna \texttt{false}
\item Saldo nunca fica negativo
\end{itemize}

\paragraph{Exercício 4 --- Programa complexo: inventário}

Crie um sistema de inventário simples com:
\begin{itemize}
\item Classe \texttt{Item} (nome, quantidade, preço unitário)
\item Classe \texttt{Inventario} (lista de itens, adicionar, remover, buscar, valor total)
\item Classe \texttt{InterfaceInventario} (menu com opções: adicionar, remover, buscar, listar, total)
\item Testes para a classe \texttt{Inventario}
\end{itemize}

\end{document}
